\section{Scope and completed results}

This blueprint records the \emph{proved} canonical Berend plan for a complex
cubic field. It includes normalized arc placement, Propositions 4.2 and 4.3,
canonical B1, coprime Proposition IV.1, the full multi-invariance dichotomy,
and Theorem 3.6. Every declaration linked by a \texttt{\textbackslash lean}
annotation is present in the current root and was compiled in the final project
build.

The statement termed an arbitrary-family corrected Berend theorem is not asserted
separately in this blueprint. The generic introduction theorems use an explicit
corrected proposition \(H\), while the canonical model constructs it. The two
independent computational inputs in \texttt{Main.lean} are outside this proof
and are not reported as mathematical theorems proved here.

\begin{definition}[The corrected unit--prime model and residual]
  \label{def:unit-prime-residual}
  \lean{CubicTwoStage.UnitPrimeModel,CubicTwoStage.NumberFieldCubic.canonicalFieldModel}
  \leanok
  Let \(K\) be a complex cubic field with chosen unit \(\varepsilon\) and
  \(\pi\in\mathcal O_K\). Write \(L\) for the canonical algebraic-integer
  lattice and \(\mathbb T\) for \(Pt/L\). The unit--prime multipliers are
  \(\mu_{u\pi^j}\), with \(u\in\mathcal O_K^\times\) and \(j\geq0\).
  Stage I is the following statement: every nonempty closed forward-invariant
  \(Y\subseteq\mathbb T\) satisfying
  \(0\in\operatorname{closure}(Y\setminus\{0\})\) equals \(\mathbb T\).
\end{definition}

\begin{lemma}[Corrected independence and expansion]
  \label{lem:unit-prime-independence}
  \lean{CubicTwoStage.BerendModulusOneArc.multIndepPair_canonicalFieldModel,CubicTwoStage.BerendUnitNoncritical.exists_mult_margin_of_pow_ne_one}
  \leanok
  Corrected multiplicative independence follows from unit/prime arithmetic.
  Infinite order of \(\varepsilon\) supplies a multiplier with real and
  complex expansion margins.
\end{lemma}
\begin{proof}
  The norm of an infinite-order unit prevents all absolute moduli from being
  \(1\).  The infinite choose--expand argument in
  \texttt{BerendUnitNoncritical.lean} supplies the margin.
  Injectivity of the complex embedding pulls a multiplicative relation back to
  the field; unit norms, the prime absolute norm, and non-torsion of the unit
  then force both exponents to vanish.
\end{proof}

\begin{lemma}[Canonical B1]
  \label{lem:canonical-b1}
  \lean{CubicTwoStage.BerendCanonicalB1.exists_canonical_B1_intLact}
  \leanok
  \uses{lem:unit-prime-independence}
  There is one of \((1,1)\) or \((\varepsilon,1)\) whose every positive power
  has irreducible characteristic polynomial over \(\mathbb Z\).
\end{lemma}
\begin{proof}
  If a prime-norm element \(\alpha\) has a positive rational power, absolute
  norms and the real embedding imply \(\alpha^3=\pm p\). Applying this to both
  \(\pi\) and \(\varepsilon\pi\) forces \(\varepsilon^6=1\), a contradiction.
  One multiplier therefore has no positive rational power. Transport through
  an integral lattice basis identifies the complexified characteristic
  polynomial with the three embeddings. Absence of rational roots and cubic
  degree give irreducibility over \(\mathbb Q\), and Gauss's lemma gives the
  integral conclusion.
\end{proof}

\section{Stage I: normalized arcs and closed ray subgroups}

\begin{definition}[Normalized samples and compact phase lattice]
  \label{def:normalized-samples}
  \lean{CubicTwoStage.BerendNormalizedCompactArc.normalizedPoint,CubicTwoStage.BerendNormalizedCompactArc.normalizedSample}
  \leanok
  For complex multipliers \(c,d\), set
  \[
    \alpha=\frac{\log\|d\|}{\log\|c\|},\qquad
    \beta=\operatorname{phaseOf}(c),\qquad \delta=\operatorname{phaseOf}(d)-\alpha\beta.
  \]
  The normalized samples are \(c^{-\lfloor j\alpha\rfloor}d^j\), and
  \(\operatorname{normalizedPoint}(\alpha,\delta)\) is the corresponding point
  of \((\mathbb R/\mathbb Z)^2\).
\end{definition}

\begin{lemma}[Non-torsion and integer directions]
  \label{lem:phase-direction}
  \lean{CubicTwoStage.BerendNormalizedCompactArc.not_isOfFinAddOrder_normalizedPoint,CubicTwoStage.BerendNormalizedCompactArc.exists_integer_direction}
  \leanok
  \uses{def:normalized-samples}
  Corrected B3 gives non-torsion of \(\operatorname{normalizedPoint}(\alpha,\delta)\).
  The lift of its compact cyclic closure contains a nonzero integral direction
  whose real line lies in that closure. In the irrational-rate case the radial
  coefficient is strictly positive.
\end{lemma}
\begin{proof}
  The log--phase exponent identity reduces a normalized rational relation to
  multiplicative dependence, forbidden by B3. The full lift is closed and
  non-discrete. The integral-functional separation theorem produces a real
  line, while integer-valuedness on the standard lattice lets it be turned into
  an integral direction. No phase-density predicate participates in this step.
\end{proof}

\begin{lemma}[Backward-orbit normalized samples]
  \label{lem:backward-orbit}
  \lean{CubicTwoStage.BerendNormalizedCompactArc.exists_backward_orbit_complex_of_acc,CubicTwoStage.BerendNormalizedCompactArc.normalizedSample_mul_basepoint_mem}
  \leanok
  \uses{def:normalized-samples}
  A closed, forward-invariant subset of \(\mathbb C\) accumulating at zero
  admits an expanding backward orbit \(c\,w_{j+1}=w_j\) inside it. Every
  normalized sample times the basepoint belongs to the set.
\end{lemma}
\begin{proof}
  First hitting times bound compact charts of bounded backward iterates.
  Compactness has a convergent subnet, whose coordinatewise limits form the
  required orbit. If \(c\neq0\), the explicit formula \(w_j=c^{-j}w_0\) makes
  membership of the normalized samples follow by forward invariance under \(d\).
\end{proof}

\begin{theorem}[Normalized arc placement]
  \label{thm:normalized-arc}
  \lean{CubicTwoStage.BerendNormalizedCompactArc.exists_normalizedArc_canonical}
  \leanok
  \uses{lem:backward-orbit,lem:phase-direction}
  With canonical unit--prime arithmetic and the corrected planar residual, a
  nonzero normalized exponential arc on a positive-length interval belongs to
  the invariant planar lift. The curve has an integer direction and the expanding
  member produced above.
\end{theorem}
\begin{proof}
  If \(\alpha\) is irrational, fractional-part convergence along the integer
  line puts the normalized spiral in the closed set. If \(\alpha=p/q\) is
  rational, the sample sequence at multiples of \(q\) reduces to a powers
  sequence of one unit-modulus phase; clearing the denominator retains
  non-torsion and yields the elliptic arc. Both proofs use actual membership of
  the backward-orbit samples, not merely density in an abstract multiplier
  group. This covers the modulus-one and non-modulus-one situations uniformly.
\end{proof}

\begin{theorem}[Exact arcThrough interface]
  \label{thm:arc-through}
  \lean{CubicTwoStage.BerendStageIAssembly.exists_arcThrough_of_normalizedArc_mem}
  \leanok
  \uses{thm:normalized-arc}
  Every normalized arc is exactly an \texttt{arcThrough} arc after positive
  rescaling of time.
\end{theorem}
\begin{proof}
  With unrestricted rate \(z\), choose \(N\) such that
  \(|\operatorname{Im}(z)|/N<\pi\). Then
  \(\log(\exp(z/N))=z/N\), so
  \(\operatorname{scaleFlow}(\exp(z/N))(Nt)=\exp(tz)\). Parameter translation
  preserves both the positive interval and each membership hypothesis; no
  approximation or weaker arc interface is used.
\end{proof}

\begin{proposition}[Infinite closed subgroup: Proposition 4.2]
  \label{prop:ray-subgroup}
  \lean{CubicTwoStage.BerendRaySubgroup.exists_translate_raySubgroup_of_arc}
  \leanok
  \uses{thm:arc-through}
  The normalized arc yields an infinite closed additive subgroup \(H\) and a
  basepoint \(b\) such that \(b+H\subseteq Y\).
\end{proposition}
\begin{proof}
  Taylor estimates transform expanding images of the arc into arbitrarily fine
  nets of arbitrarily long segments. Compactness gives a limiting ray. Closure
  of a ray is an additive subgroup: a compact subsequence produces return times,
  and predecessor times converge to inverses. Preconnectedness and discreteness
  of the lattice give infinitude. The basepoint of the ray need not be planar,
  and \(H\) is not asserted to be invariant.
\end{proof}

\section{Proposition 4.3: characters, van der Waerden, and nets}

\begin{lemma}[Strict rank drop of the annihilator]
  \label{lem:annihilator-rank}
  \lean{CubicTwoStage.BerendProp43Canonical.CanonicalProp43.finrank_rationalAnnihilatorSpan_lt}
  \leanok
  The rational span of the character annihilator of an infinite subgroup has
  dimension strictly below the ambient character-space dimension.
\end{lemma}
\begin{proof}
  If the annihilator has finite index, coordinate characters restrict the
  subgroup to a finite product of finite coordinate torsion sets. This makes
  the subgroup finite. Hence an infinite subgroup has infinite annihilator
  index. Rationalization is injective and carries this to the required strict
  finite-rank inequality.
\end{proof}

\begin{lemma}[No finite character cover]
  \label{lem:no-character-cover}
  \lean{CubicTwoStage.BerendProp43Canonical.CanonicalProp43.not_finite_cover_latticeCharacter_annihilator}
  \leanok
  \uses{lem:canonical-b1,lem:annihilator-rank}
  No finite family of nonzero lattice characters can have one of its members
  annihilate every positive iterate of an infinite subgroup.
\end{lemma}
\begin{proof}
  Apply van der Waerden to obtain a monochromatic arithmetic progression of
  length equal to the ambient dimension. Shift to its initial point. B1 makes
  the first dimension-many covector iterates linearly independent: a relation
  would be a nonzero polynomial of smaller degree whose kernel is invariant,
  contradicting irreducibility. They therefore cannot all lie in the
  lower-dimensional annihilator span.
\end{proof}

\begin{theorem}[Proposition 4.3 in net form]
  \label{thm:unit-prime-net}
  \lean{CubicTwoStage.BerendProp43Canonical.CanonicalProp43.exists_canonical_unitPrime_net}
  \leanok
  \uses{lem:no-character-cover}
  For canonical B1, each infinite closed \(H\) and each \(\delta>0\), some
  iterate of \(H\) is a \(\delta\)-net of the entire torus.
\end{theorem}
\begin{proof}
  Finite character prefixes and arbitrarily late good times give a strictly
  increasing sequence on which every nonzero character is eventually nontrivial.
  Normalized Haar measures on the compact image groups have vanishing nontrivial
  character integrals. Stone--Weierstrass identifies approximations for all
  continuous functions. A positive bump-supported ball argument converts this
  into eventual intersection of every ball. Finitely many such balls cover the
  compact torus, delivering one iterate that is a net. No invariance of \(H\)
  is required.
\end{proof}

\begin{lemma}[Translated-subgroup last mile]
  \label{lem:translate-last-mile}
  \lean{CubicTwoStage.BerendStageIProp43.eq_univ_of_iterate_net_translate}
  \leanok
  \uses{thm:unit-prime-net}
  If \(Y\) is closed and forward-invariant and contains an iterate of \(b+H\),
  the subgroup nets force \(Y\) to be the whole torus.
\end{lemma}
\begin{proof}
  Center the net at \(x-\sigma^n b\). Translation by \(\sigma^n b\) produces a
  point of the iterate of \(b+H\) within \(\delta\) of \(x\). Forward invariance
  keeps that iterate in \(Y\), so every ball around \(x\) meets \(Y\). Closedness
  yields \(x\in Y\).
\end{proof}

\section{Stage I, full dichotomy, and Theorem 3.6}

\begin{theorem}[Stage I from a normalized arc]
  \label{thm:stage-one-arc}
  \lean{CubicTwoStage.BerendStageIProp43.eq_univ_canonical_of_arcThrough}
  \leanok
  \uses{thm:arc-through,prop:ray-subgroup,lem:translate-last-mile}
  An exact normalized arc in \(Y\) proves \(Y=\mathbb T\).
\end{theorem}
\begin{proof}
  Use the tangent net and Proposition 4.2 to obtain \(b+H\), then the B1/net
  endomorphism and the translated-subgroup last mile.
\end{proof}

\begin{lemma}[Off-plane versus on-plane accumulation]
  \label{lem:stage-one-split}
  \lean{CubicTwoStage.BerendStageIPlanarOnly.zeroNonisolatedFull_canonical_planarOnly'}
  \leanok
  \uses{lem:unit-prime-independence}
  Canonical Stage I reduces to the corrected planar residual, without an
  expansive-direction hypothesis.
\end{lemma}
\begin{proof}
  Real multiplier density and coordinate-axis density handle off-plane
  accumulation. Non-criticality supplies \(|\sigma_1\varepsilon|\neq1\), and
  replacing \(\varepsilon\) by \(\varepsilon^{-1}\) if necessary handles either
  expansion direction.
\end{proof}

\begin{theorem}[Completed canonical Stage I]
  \label{thm:canonical-stage-one}
  \lean{CubicTwoStage.BerendStageIAssembly.zeroNonisolatedFull_canonical}
  \leanok
  \uses{thm:normalized-arc,thm:arc-through,thm:stage-one-arc,lem:stage-one-split}
  Every nonempty closed forward-invariant subset of the canonical torus that has
  zero as a non-isolated point is the entire torus.
\end{theorem}
\begin{proof}
  Apply the off-plane/on-plane split. On the complex plane, transport the
  corrected residual to nonzero complex accumulation and use the normalized arc
  and Propositions 4.2 and 4.3 established above.
\end{proof}

\begin{theorem}[Coprime Proposition IV.1]
  \label{thm:coprime-iv1}
  \lean{CubicTwoStage.BerendIV1Arith.propIV1_unitPrime'_canonical}
  \leanok
  Let \(p,q\) be distinct rational primes with \(\lvert\operatorname{Nm}\pi\rvert=p\).
  Coprime Proposition IV.1 holds for the canonical unit--prime family.
\end{theorem}
\begin{proof}
  Ideal norms supply coprimality and degree-one quotient representatives after
  saturation on multiplicative levels. The unit--prime lift preserves these
  representatives under coprime levels and gives the injectivity needed by
  Proposition IV.1.
\end{proof}

\begin{theorem}[Full canonical multi-invariance]
  \label{thm:canonical-dichotomy}
  \lean{CubicTwoStage.BerendStageIAssembly.dichotomy_canonical}
  \leanok
  \uses{thm:canonical-stage-one,thm:coprime-iv1}
  Every closed nonempty forward-invariant proper subset of the canonical torus is finite.
\end{theorem}
\begin{proof}
  The expanding margin supplies finite kernels and finite torsion. Stage I,
  Proposition IV.1, and the §5 propagation assembly therefore give the full
  multi-invariance dichotomy.
\end{proof}

\begin{theorem}[Canonical Theorem 3.6]
  \label{thm:canonical-uncovered}
  \lean{CubicTwoStage.NumberFieldCubic.thm_uncovered_canonical}
  \leanok
  \uses{thm:canonical-dichotomy,thm:coprime-iv1}
  If \(p\) is prime and the usual unit/prime, \(p\)-adic, B1, and B2 hypotheses hold,
  the uncovered set lies in finitely many lattice classes, each having a field
  representative.
\end{theorem}
\begin{proof}
  Choose \(q\) prime and larger than \(p\). The canonical dichotomy and coprime
  IV.1 construct the corrected proposition \(H\). Closure of the uncovered
  orbit is finite by \(H\); injectivity of the unit action gives periodic
  unit orbits, and solving the denominator equation produces field
  representatives. The rational primality needed by IV.1 is explicit; a
  composite absolute norm is not silently treated as prime.
\end{proof}

\begin{corollary}[Generic corrected paper statements]
  \label{cor:generic-paper}
  \lean{CubicTwoStage.thm_uncovered,CubicTwoStage.thm_uncovered_field,CubicTwoStage.thm_intro_equiv}
  \leanok
  \uses{thm:canonical-dichotomy}
  The same closure-of-orbits, periodic-orbit, representative, and geometric
  equivalence arguments prove the generic uncovered and introduction theorems
  using an explicit corrected proposition \(H\).
\end{corollary}
\begin{proof}
  This is precisely the corrected wrapper interface: arbitrary models accept
  \(H\) as a hypothesis, while canonical models construct it using
  Theorem~\ref{thm:canonical-dichotomy}.
\end{proof}

\section{Trust boundary and verification}

Every Berend-facing declaration above audits with only \texttt{propext},
\texttt{Classical.choice}, and \texttt{Quot.sound}. The old false literal axiom
was removed together with its legacy consumers. No \texttt{sorry} or substitute
axiom is used.

The independent \texttt{thm\_main\_computational\_input} and
\texttt{thm\_smooth\_computational\_input} remain only as the separately
identified computational data in \texttt{Main.lean}; neither enters this proof.
The complete project was validated by the zero-exit full build
\texttt{jmw9uaao4}. The labels in this file supply the plasTeX dependency graph
and the accompanying Graphviz graph.
